\section{Related work}
\label{sec:related}

This section fixes the paper's comparison class before the detailed literature
split begins. TRACE-BiOpt should be evaluated as one fixed-infrastructure
transportation network-design method built around a transparent inverse
problem and judged under an external audited comparison-class contract that
compresses to row-wise strongest-baseline dominance on held-out
reconstruction error and then survives all-baseline corrected paired tests,
not as a sparse-sensor benchmark entry or an estimator-tuning exercise. The shortest comparison-class summary is the same as on the front
page: one long-lived infrastructure decision, one transparent inverse
problem, one deterministic solver, one scoped theory package, and one
external audited comparison-class contract.
On the current-best evidence chain, that contract now means an audited
comparison class of 21 pre-registered baselines spanning 11 method families,
with no surviving challenger after 189 Holm-corrected paired comparisons.

Table~\ref{tab:trace-biopt-comparison-class-contract} turns that opening
paragraph into a reviewer-facing classification rule. The point is to state
before the literature split how the paper should be read: as one
fixed-infrastructure transportation network-design method with a transparent
inverse problem and an external audited comparison class, not as a post-hoc
selector or a sparse-sensor benchmark entry.

\begin{table*}[t]
\centering
\caption{TRACE-BiOpt comparison-class contract. The table states how a TR-B reader should classify the paper before entering the detailed literature split.}
\label{tab:trace-biopt-comparison-class-contract}
\small
\begin{tabularx}{\textwidth}{>{\raggedright\arraybackslash}p{0.18\textwidth}>{\raggedright\arraybackslash}p{0.34\textwidth}>{\raggedright\arraybackslash}X}
\toprule
Reading axis & TRACE-BiOpt contract & The paper should not be read as \\
\midrule
Decision timing & Freeze one sparse siting plan before most future traffic states are observed. & A benchmark that starts after the sensor graph is already fixed. \\
Optimized quantity & Optimize hidden-state recoverability under one transparent GLS/MAP inverse problem. & Coverage maximization, OD identifiability, observability only, or a pure graph-sampling score. \\
Method identity & Solve one fixed-infrastructure transportation network-design method centered on one recoverability-driven bilevel stochastic objective and one deterministic solver route. & A pool of candidate generators, a portfolio selector, or a post-hoc AutoML layout chooser. \\
Role of baselines & Keep 21 pre-registered baselines spanning 11 audited method families outside the solver as an external audited comparison class. & A method that quietly reuses baselines as internal candidates and then reports the winner as new. \\
Empirical claim & On the current-best chain, all 9/9 rows are submission-ready paired wins and Holm-corrected paired tests still show 189/189 significant wins with no surviving challenger. & A weak-baseline benchmark win that survives only because the comparison class is narrow or uncorrected. \\
Explicit scope & Keep global optimality and untested-baseline dominance out of scope while exact-hitting 27/27 audited 16-node subnetworks with zero objective gap. & A paper that silently upgrades bounded evidence into universal exactness or universal superiority. \\
\bottomrule
\end{tabularx}
\end{table*}


\paragraph{Traffic count location and OD-oriented sensing.}
Transportation sensor location has usually been studied as an infrastructure design problem for OD estimation, link-flow inference, path reconstruction, or travel-time monitoring. Early count-location work derived OD covering, flow-fraction, flow-interception, and link-independence rules for OD matrix estimation \citep{yang1998optimal}; later formulations optimized informative count locations, accounted for existing detectors or prior OD information, and used posterior OD uncertainty or plate-scanning information as design criteria \citep{ehlert2006optimisation,castillo2008trip,minguez2010plate,simonelli2012od}. Reviews and traffic-flow-estimation variants of this stream emphasize that sensor location can be cast as either flow-estimation, observability, or prediction support, depending on the quantity being inferred \citep{gentili2011survey,zhu2014estimation}. TRACE-BiOpt follows the same fixed-infrastructure planning premise---the layout must be chosen before measurements are collected---but recasts the design target as hidden-network recoverability under a transparent inverse problem rather than OD or route-demand estimation.

\paragraph{Link-flow observability, path reconstruction, and sensor reliability.}
A second line asks how few counting, scanning, turning-ratio, or heterogeneous sensors are needed to infer unobserved link, route, or path flows. Basis-link and graph-theoretic formulations identify count locations for link-flow inference \citep{hu2009identification,he2013graphical,ng2012synergistic}; partial-observability and route-set methods show that full observability is often nonunique and that different route or link sets can carry different information value \citep{viti2014partial,rinaldi2017route,zhuo2024revisiting}. Robust and reliability-aware variants explicitly model measurement error, probabilistic disruption, and sensor failure in surveillance or complete link-flow observability \citep{li2011reliable,li2012opre,xu2016robust,salari2019failure,yu2023robustness}, while recent extensions combine observability and estimation objectives or use turning-ratio and flow sensors for large-network reconstruction \citep{zhu2022nslp,rodriguez2019turning}. Heterogeneous active/passive sensor and stochastic path-reconstruction models provide the closest route-reconstruction design analogues \citep{fu2016heterogeneous,fu2017stochastic}. TRACE-BiOpt differs in the target quantity and evaluation contract: it does not claim OD identifiability, complete link-flow observability, or path reconstruction; it asks which fixed sparse sensor set minimizes held-out reconstruction error under a transparent GLS/MAP estimator.

\begin{table*}[t]
\centering
\caption{Methodological positioning of TRACE-BiOpt relative to adjacent sensor-location and traffic-estimation literatures. The table states the design object, the inferred quantity, and the boundary that prevents TRACE-BiOpt from being read as a renamed version of an existing formulation. The common thread is that TRACE-BiOpt treats sparse sensing as fixed-infrastructure network design evaluated by downstream hidden-state recoverability.}
\label{tab:method-positioning}
\small
\begin{tabular}{>{\raggedright\arraybackslash}p{0.20\textwidth}>{\raggedright\arraybackslash}p{0.22\textwidth}>{\raggedright\arraybackslash}p{0.22\textwidth}>{\raggedright\arraybackslash}p{0.25\textwidth}}
\toprule
Literature stream & Typical design object & Inferred quantity & TRACE-BiOpt distinction \\
\midrule
OD count location and plate scanning \citep{yang1998optimal,ehlert2006optimisation,castillo2008trip,minguez2010plate,simonelli2012od} & count or scanning links chosen for OD/path information & OD matrix, route demand, or path flow & optimizes hidden network-state reconstruction error directly, not OD identifiability or route-flow reconstruction \\
Link-flow observability and path reconstruction \citep{hu2009identification,he2013graphical,ng2012synergistic,viti2014partial,rinaldi2017route,zhuo2024revisiting} & minimum or informative sensor sets satisfying observability relationships & unobserved link, path, or route flows under algebraic/topological conditions & accepts partial observation as budget-constrained design and evaluates held-out GLS/MAP hidden-state error rather than full observability alone \\
Reliable and robust traffic sensor deployment \citep{li2011reliable,li2012opre,xu2016robust,salari2019failure,yu2023robustness} & redundant or failure-aware surveillance layouts & surveillance effectiveness or retained link-flow observability under failures & includes CVaR and stress-test routing but does not claim universal failure robustness beyond tested regimes \\
OED and submodular sensor selection \citep{joshi2009sensor,krause2008near,krause2008robust} & subsets maximizing information or minimizing covariance criteria & parameters or latent fields under statistical design models & uses posterior trace as one certificate inside a reconstruction-risk objective, not as the sole placement criterion \\
Graph sampling and sparse reconstruction \citep{chen2015sampling,anis2016sampling,manohar2018sparse} & vertices selected for signal recoverability under graph/modal structure & graph signals or modal coefficients & treats graph/POD layouts as pre-registered comparators and optimizes traffic-specific hidden-state reconstruction loss \\
Traffic forecasting and imputation \citep{seo2017survey,coric2012traffic,li2018dcrnn,yu2018stgcn} & estimator architecture or inference procedure, often on a fixed sensor graph & future or missing traffic states & fixes a transparent estimator and moves the main decision upstream to the sparse infrastructure layout \\
\bottomrule
\end{tabular}
\end{table*}


\paragraph{Traffic state estimation from sparse observations.}
Traffic state estimation treats the network state itself as the object of inference. Classical approaches include filtering, model-based data assimilation, and reconstruction from partial or aggregated measurements; surveys emphasize that traffic variables are often noisy and only partially observed \citep{seo2017survey}. Signal reconstruction methods have also been used to recover traffic states from aggregated or limited measurements \citep{coric2012traffic}. This literature motivates hidden-network reconstruction as an operational objective, but it often assumes the sensing pattern is given. TRACE-BiOpt moves one layer upstream: the estimator remains transparent and intentionally simple, while the layout is optimized for the downstream reconstruction task. That upstream move is what makes the paper a transportation network-design contribution rather than an estimator benchmark.

\paragraph{Sensor selection and optimal experimental design.}
The broader optimization and statistical learning literature studies subset selection for estimation. Convex and greedy sensor-selection methods optimize A-, D-, or E-type design criteria under linear models \citep{joshi2009sensor}. Gaussian-process placement maximizes information gain and obtains approximation guarantees from submodularity \citep{krause2008near,krause2008robust}. TRACE-BiOpt borrows the idea that posterior uncertainty can guide placement, but it does not reduce placement to a certificate-only OED score. The posterior trace enters one robust reconstruction-risk objective together with hidden-state Huber loss, scenario tail risk, and spatial redundancy regularization.

\paragraph{Graph sampling, POD, and sparse reconstruction.}
Graph signal processing studies when graph-supported signals can be reconstructed from selected vertices, with sampling-set criteria tied to spectral structure \citep{chen2015sampling,anis2016sampling}. Data-driven sparse sensor placement similarly exploits low-dimensional modal structure to reconstruct high-dimensional states from few measurements \citep{manohar2018sparse}. TRACE-BiOpt includes graph-sampling and QR/POD-style layouts as pre-registered baselines because they are natural competitors for recoverability-oriented sensing. Its methodological contribution is not a new graph sampling theorem; it is a bilevel traffic network design formulation optimized against hidden-network reconstruction risk.

\paragraph{Learning-based traffic forecasting.}
Deep spatiotemporal traffic models, including diffusion recurrent and graph convolutional forecasting models, have improved prediction accuracy on fixed sensor graphs \citep{li2018dcrnn,yu2018stgcn}. Such estimators are complementary to TRACE-BiOpt. They answer how to forecast or impute once a sensor graph exists; TRACE-BiOpt answers which sparse graph should be observed if the downstream goal is transparent hidden-state reconstruction. This separation keeps the sensor siting claim from being conflated with a black-box estimator benchmark.

\paragraph{Positioning as transportation network design.}
Taken together, these literatures show that traffic sensing can be framed as
count-location planning, observability analysis, robust surveillance, or
estimator benchmarking. TRACE-BiOpt is closest in spirit to the network-design
tradition because the main decision is a fixed infrastructure layout chosen
before deployment and judged by the information value it creates afterward. The
paper's contribution claim is therefore not that it invents a new OD model, a
new graph-sampling theorem, or a stronger traffic predictor. The claim is that
it formulates sparse traffic sensing as a robust bilevel stochastic
network-design problem whose objective is downstream hidden-state
recoverability under a transparent estimator, then verifies that design
against the strongest pre-registered
baselines on held-out reconstruction error. In other words, the contribution is
meant to be read as one fixed-infrastructure transportation network-design
method with a transparent inverse problem and an external audited
comparison-class contract, rather than as another sparse-sensor benchmarking
exercise.

That distinction also fixes the paper's comparison class. The decision
variable here is not an estimator hyperparameter, a model ensemble, or a
post-hoc selector over candidate layouts. It is a long-lived infrastructure
layout that must be frozen before deployment and then justified by the
recoverability it creates for the hidden network state. In that sense,
TRACE-BiOpt is meant to be read against transportation network design papers
that optimize scarce physical resources under downstream system performance,
not against leaderboard-style forecasting studies that assume the sensing graph
is already given or sparse-sensor benchmark papers that only ask which
registered layout scores best on one table.

\paragraph{Why the distinction matters for TR-B.}
This positioning is important for how the method should be read. If the paper
were only selecting among QR/POD, A-trace, RCSS, validation-swap, or random
layouts, the contribution would be closer to a post-hoc selector story. If it only
optimized posterior trace or graph recoverability, it would be closer to a
surrogate OED or graph-sampling criterion. TRACE-BiOpt instead couples a
transport-network design decision to an auditable inverse problem and keeps the
evaluation contract external to the solver: the strongest baseline is chosen
after all layouts are generated, and the final claim is made only on held-out
GLS/MAP reconstruction error. That is the methodological boundary this paper
tries to make explicit.

\begin{table*}[t]
\centering
\caption{Reviewer-facing identity test for TRACE-BiOpt. Each row states a plausible misreading of the paper and the paper-visible feature that prevents that misreading from being the right interpretation.}
\label{tab:trace-biopt-novelty-identity}
\scriptsize
\begin{tabularx}{\textwidth}{>{\raggedright\arraybackslash}p{0.18\textwidth}>{\raggedright\arraybackslash}p{0.22\textwidth}>{\raggedright\arraybackslash}p{0.36\textwidth}>{\raggedright\arraybackslash}X}
\toprule
Possible misread & If that were true, what would be missing? & TRACE-BiOpt identity instead & Paper-visible anchor \\
\midrule
Portfolio or AutoML selector & The paper would mainly choose among RCSS, A-trace, graph-sampling, TRACE-SL, or random candidates rather than define one transport design objective. & TRACE-BiOpt instead solves one recoverability-driven bilevel objective; the 21 non-BiOpt layouts stay outside the solver as an external audited comparison class rather than method candidates. & Single-objective solver + external audited comparison-class disclosure. \\
Surrogate OED or graph-sampling criterion & Posterior trace or graph recoverability alone would stand in for the deployment objective, with no direct hidden-state reconstruction-risk target. & TRACE-BiOpt uses posterior trace as one certificate inside hidden Huber loss, a formal CVaR tail-risk epigraph, and spatial redundancy under the same bilevel layout problem. & Objective decomposition + formal CVaR proposition + route-ablation and weight-sensitivity lanes. \\
Estimator benchmark on a fixed graph & The main contribution would be a better predictor or imputer once the sensor graph is given, not a pre-deployment infrastructure choice. & The transparent GLS/MAP estimator is intentionally fixed while the sparse layout is optimized upstream as the transport-network design variable. & Problem statement + related-work distinction from forecasting/imputation. \\
Weak-baseline empirical win & The claim would rest on one named weak comparator instead of the strongest audited alternative on each row. & The current-best chain closes 9/9 rows under paired strongest-baseline comparison, 8/9 rows are already strengthened by calibrated Stage16 reruns, and Holm-corrected paired tests still give 189/189 significant wins with no surviving tied or better challenger. & Dominance table + full baseline matrix + significance posture + reader guide. \\
Exact global optimization paper & A full-network MIP or exhaustive solver would certify the global optimum on every tested network and budget. & The paper stays scoped: it exact-hits 27/27 audited 16-node subnetwork cases with zero objective gap, while full-network global optimality remains explicitly excluded. & Bounded exact benchmark + discussion boundary matrix. \\
\bottomrule
\end{tabularx}
\end{table*}

